\section{From an experiment to a sample space}

The previous chapter introduced the formal language of probability: a sample
space, events as subsets of the sample space, operations on events, and a
probability assignment. We now use that language in concrete models. The main
point of the present chapter is methodological:
\[
\text{a sample space is not found inside the experiment; it is constructed.}
\]

A physical or conceptual experiment may be described in ordinary language, but
ordinary language does not by itself say what the elementary outcomes are. We
must decide what information is recorded, what distinctions are visible in the
model, and what distinctions are deliberately ignored.

\begin{definitionbox}
A \textbf{probabilistic model} begins with a decision about what is recorded. A
\textbf{sample space} \(\Omega\) is the set of all elementary outcomes under
that recording rule. A \textbf{probability assignment} then gives probabilities
to events, that is, to subsets of \(\Omega\).
\end{definitionbox}

Thus the construction has the following order:
\[
\text{experiment}
\longrightarrow
\text{recording rule}
\longrightarrow
\text{elementary outcome}
\longrightarrow
\Omega
\longrightarrow
\text{events}
\longrightarrow
P.
\]

The order matters. If we skip the recording rule, then we do not really know
what one elementary outcome is. If we do not know what one elementary outcome is,
then we do not know what the sample space is. If we do not know the sample space,
then symbols such as \(A\subseteq\Omega\) and \(P(A)\) have no precise meaning.

\subsection*{The same physical situation may have different sample spaces}

Consider the instruction:
\[
\text{toss a coin twice.}
\]
This instruction describes a physical procedure, but it does not yet fully
describe the mathematical model. Several recording rules are possible.

If we record the full ordered sequence of outcomes, then one elementary outcome
is a sequence such as \(HT\). The sample space is
\[
\Omega_1=\{HH,HT,TH,TT\}.
\]
In this model \(HT\neq TH\), because the first position and the second position
are different pieces of information.

If instead we record only the number of heads, then one elementary outcome is a
number. The sample space is
\[
\Omega_2=\{0,1,2\}.
\]
In this model the sequences \(HT\) and \(TH\) are not separate elementary
outcomes. They both correspond to the same recorded value, namely \(1\).

If we record only whether at least one head appeared, then the sample space is
\[
\Omega_3=\{\text{yes},\text{no}\}.
\]
This is another legitimate model, but it answers a different question.

\begin{remarkbox}
The phrase ``the experiment is the same'' can be misleading. The physical action
may be the same, but the mathematical sample space depends on what is recorded.
Probability theory does not work directly with the physical action; it works
with the mathematical description produced from that action.
\end{remarkbox}

\subsection*{Elementary outcomes are model-dependent}

An elementary outcome is not necessarily the smallest physical object involved.
It is the smallest complete result according to the model.

\begin{examplebox}
Suppose a ball is drawn from an urn containing three red balls and two blue
balls. If the balls are labelled
\[
R_1,R_2,R_3,B_1,B_2,
\]
and the label is recorded, then the sample space is
\[
\Omega=\{R_1,R_2,R_3,B_1,B_2\}.
\]
Here the elementary outcomes are individual labelled balls.

If only the colour is recorded, then the sample space is instead
\[
\Omega'=\{R,B\}.
\]
Here the elementary outcomes are colours, not labelled balls. The two sample
spaces describe different levels of information.
\end{examplebox}

The second model is not wrong. It may be exactly the model we need if the
question is only about colour. But it must be treated honestly: changing the
recording rule changes the sample space and may also change how probabilities
are assigned.

\subsection*{The sample space is only half of the model}

After \(\Omega\) has been constructed, we still have to assign probabilities.
The sample space tells us which elementary outcomes exist in the model. It does
not automatically say how likely they are.

\begin{examplebox}
For a single coin toss we may choose
\[
\Omega=\{H,T\}.
\]
If the coin is modelled as symmetric, then we assign
\[
P(H)=P(T)=\frac12.
\]
If the coin is modelled as biased toward heads, then we may assign
\[
P(H)=p,\qquad P(T)=1-p,\qquad p\neq \frac12.
\]
The sample space is the same in both models, but the probability assignment is
different.
\end{examplebox}

This distinction is fundamental. A list of possible outcomes is not yet a
probability model. A probability model requires both:
\[
\text{possible outcomes} \quad \text{and} \quad \text{probabilities of events}.
\]

\subsection*{Uniform and non-uniform assignments}

In a finite sample space
\[
\Omega=\{\omega_1,\omega_2,\ldots,\omega_n\},
\]
a probability assignment is determined by the probabilities of the elementary
outcomes:
\[
P(\{\omega_1\})=p_1,\quad P(\{\omega_2\})=p_2,\quad \ldots,\quad
P(\{\omega_n\})=p_n.
\]
These numbers must satisfy
\[
p_i\geq 0\quad\text{for every }i,
\]
and
\[
p_1+p_2+\cdots+p_n=1.
\]
Then the probability of an event is obtained by adding the probabilities of the
elementary outcomes contained in that event:
\[
P(A)=\sum_{\omega_i\in A}p_i.
\]

\begin{definitionbox}
A finite probability model is \textbf{uniform} if all elementary outcomes have
the same probability. If \(|\Omega|=n\), this means
\[
P(\{\omega\})=\frac1n
\]
for every \(\omega\in\Omega\).
\end{definitionbox}

Only in a uniform finite model do we obtain the classical formula
\[
P(A)=\frac{|A|}{|\Omega|}.
\]
This formula is not the definition of probability in every finite model. It is a
consequence of the special assumption that all elementary outcomes are equally
likely.

\begin{remarkbox}
A common error is to count how many named outcomes appear in a sample space and
then assume that all names have equal probability. Equal naming is not the same
as equal probability. Equal probability must come from the mechanism of the
model or from an explicit modelling assumption.
\end{remarkbox}

\subsection*{Events are questions asked of the sample space}

Once \(\Omega\) is fixed, an event is a subset of \(\Omega\). Conceptually, an
event answers a question about the elementary outcome.

\begin{examplebox}
If a die is rolled and
\[
\Omega=\{1,2,3,4,5,6\},
\]
then the question
\[
\text{is the result even?}
\]
selects the event
\[
A=\{2,4,6\}.
\]
The question
\[
\text{is the result at least }5?
\]
selects the event
\[
B=\{5,6\}.
\]
The compound question
\[
\text{is the result even and at least }5?
\]
selects the event
\[
A\cap B=\{6\}.
\]
\end{examplebox}

This is why the sample space must be fixed before events are discussed. The same
sentence can select different sets in different spaces.

\subsection*{A first complete modelling pattern}

When solving elementary probability problems, we shall repeatedly use the
following pattern.

\begin{enumerate}
    \item Describe the physical or conceptual experiment.
    \item State what information is recorded.
    \item Identify one elementary outcome.
    \item Construct the sample space \(\Omega\).
    \item Describe the event or events of interest as subsets of \(\Omega\).
    \item Decide whether elementary outcomes are equally likely.
    \item Assign probabilities to elementary outcomes or directly to events.
    \item Compute the required probability and interpret it in the original situation.
\end{enumerate}

This pattern may look slow. That is intentional. At the beginning, speed is less
important than correctness. Most wrong solutions to elementary probability
problems do not fail because of difficult arithmetic. They fail because the
sample space, the elementary outcomes, or the probability assignment were never
made explicit.

\begin{remarkbox}
The goal of this chapter is not to collect examples. The goal is to learn a
method: construct the model first, calculate only after the model is clear.
\end{remarkbox}
